\section{Preliminaries}\label{framework}
We start with some notation for fractional tilings and covers.

Throughout this section, $H$ is an orientation of a cycle of order $r\ge4$.
For a homomorphism
$f$ from $H$ to $F$, recall that $m_f(v)=|f^{-1}(v)|$. A \emph{fractional
homomorphic $H$-tiling} of $F$ is a family of weights $\lambda_f\ge0$ on the homomorphisms $f$ from $H$ to $F$ satisfying
$\sum_fm_f(v)\lambda_f\le1$ for every $v\in V(F)$, and it is \emph{perfect} if
\begin{equation}\label{framework390:loads}
 \sum_f m_f(v)\lambda_f=1\quad\text{for every }v\in V(F).
\end{equation}
Recall that a fractional $H$-cover of $F$ is a function $w:V(F)\to\mathbb R_{\ge0}$
such that every homomorphism from $H$ to $F$ has cost at least one.
A homomorphism is \emph{tight} with respect to $w$ if its cost equals one.
We say that a fractional $H$-cover of $F$ is \emph{deficient} if its total weight is less than $|F|/r$; by Lemma~\ref{general:LP} below, Theorem~\ref{compact390:finite} is equivalent to the statement that no fractional $H$-cover of $F$ is deficient.

\begin{lemma}[Linear programming duality]\label{general:LP}
A digraph $F$ of order $N$ has a perfect fractional homomorphic
$H$-tiling if and only if every fractional $H$-cover of $F$ has total weight
at least $N/r$. Moreover, the following hold.
\begin{props}
\item If $F$ has a perfect fractional homomorphic $H$-tiling, then it has one in which at most $N$ homomorphisms have positive weight.\label{LP:support}
\item Let $w$ be a fractional $H$-cover of minimum total weight, and let $\lambda$ be a fractional homomorphic $H$-tiling of maximum total weight. Then every vertex of positive weight is covered by $\lambda$ with total multiplicity one, and every homomorphism of positive weight in $\lambda$ is tight.\label{LP:slack}
\item Every deficient fractional $H$-cover $w$ of $F$ of minimum total weight satisfies $\min w=0$ and $\max w\le1$.\label{LP:normal}
\end{props}
\end{lemma}
\begin{proof}
Consider the linear programme which maximises $\sum_f\lambda_f$ over all fractional homomorphic $H$-tilings. Its dual programme is to minimise the total weight of a fractional $H$-cover.
Since $\sum_vm_f(v)=r$ for every $f$, every fractional tiling has total weight at most $N/r$, with equality
exactly when it is perfect. Thus, the first assertion follows from the duality theorem of linear programming.

For~\ref{LP:slack}, write $l(v)=\sum_fm_f(v)\lambda_f$ and let $W$ be the total weight of $w$, so that $W=\sum_f\lambda_f$ by duality. Then
\[
 0=W-\sum_f\lambda_f
 =\sum_vw(v)(1-l(v))+
   \sum_f\lambda_f\left(\sum_{x\in V(H)}w(f(x))-1\right),
\]
and each summand on the right is nonnegative. Hence $w(v)>0$ implies $l(v)=1$, and
$\lambda_f>0$ implies that $f$ is tight.

For~\ref{LP:normal}, let $w$ be an optimal cover, with total weight $W<N/r$. If $w(v)>1$ for some $v$, then replacing $w(v)$ by one preserves the cover constraints and decreases the total weight, which is impossible. Thus $w\le1$.
If $\mu=\min w>0$, then $\mu\le W/N<1/r$, and
$(w-\mu)/(1-r\mu)$ is a cover of total weight
$(W-N\mu)/(1-r\mu)<W$, a contradiction. Thus $\min w=0$.

Finally, for~\ref{LP:support}, take a perfect fractional tiling, and suppose that more than $N$ homomorphisms have positive weight. Then their vectors $(m_f(v))_{v\in V(F)}$ are linearly dependent, say $\sum_fa_fm_f=0$ with not all $a_f$ zero. Summing over $v$ gives $r\sum_fa_f=0$, so some coefficients $a_f$ are positive and some are negative.
Let $t=\min\{\lambda_f/a_f:a_f>0\}$, and subtract $ta_f$ from $\lambda_f$ for every $f$. This preserves~\eqref{framework390:loads} and nonnegativity, and makes at least one more weight zero. Repeating this gives a perfect fractional tiling with at most $N$ homomorphisms of positive weight.
\end{proof}

If $H$ can be mapped onto another digraph $K$ so that every vertex of $K$ receives the same number of vertices of $H$ on average, then tilings by $K$ give tilings by $H$, as follows.

\begin{lemma}\label{transfer:fibres}
Let $H$ and $K$ have orders $r$ and $q$. Suppose that there is a random homomorphism
$\psi$ from $H$ to $K$ such that $\mathbb E|\psi^{-1}(x)|=r/q$ for every vertex $x$ of $K$.
Then, every perfect fractional homomorphic $K$-tiling of a digraph $F$ gives a
perfect fractional homomorphic $H$-tiling of $F$.
\end{lemma}
\begin{proof}
For each homomorphism $\phi$ from $K$ to $F$ with weight $\lambda_\phi$, and each value $\psi_0$ of $\psi$, which has probability $\theta_{\psi_0}$ say, give the homomorphism $\phi\circ\psi_0$ from $H$ to $F$ the weight
$(q/r)\lambda_\phi\theta_{\psi_0}$, adding the weights of coincident homomorphisms.
The total multiplicity of a vertex $v$ of $F$ in the resulting tiling is
\[
 \frac qr\sum_\phi\lambda_\phi
 \sum_{x\in\phi^{-1}(v)}\mathbb E|\psi^{-1}(x)|
 =\sum_\phi\lambda_\phi|\phi^{-1}(v)|=1,
\]
as required.
\end{proof}

\subsection{Partial homomorphisms and completion sets}\label{sec:prelim-completion}

Let $D$ be a digraph. A \emph{partial homomorphism} into $D$ is a homomorphism $Q$ from $H-j$ to $D$, where the deleted position $j$ of $H$ is specified as part of $Q$.
We call the two positions of $H$ adjacent to $j$ the \emph{endpoints} of $Q$.
The edges of $H$ between $j$ and its two endpoints prescribe, for each endpoint, whether the image of $j$ must be an outneighbour or an inneighbour of the image of that endpoint.
We call these two neighbourhoods the \emph{completion sets} $X_Q$ and $Y_Q$ of $Q$, and $(X_Q,Y_Q)$ the \emph{completion pair} of $Q$.
Thus, $Q$ extends to a homomorphism of $H$ precisely by mapping $j$ to a vertex of $X_Q\cap Y_Q$.
Given weights on $V(D)$, a partial homomorphism is \emph{minimum} if its cost is minimum among all partial homomorphisms into $D$, over all deleted positions.

We will often consider a finite measure $\rho$ on partial homomorphisms, whose total mass we call its \emph{mass}. Its \emph{multiplicity function} is
\[
 p(v)=\sum_Q\rho(Q)\,|Q^{-1}(v)|,
\]
that is, $p(v)$ is the total weight, under $\rho$, of the positions mapped to $v$. As each partial homomorphism has $r-1$ positions, a measure of mass $m$ has total multiplicity $p(V(D))=(r-1)m$. For a partial homomorphism $Q$, we write
\[
 \Phi_p(Q)=p(X_Q)+p(Y_Q),
\]
and we say that the completion pair of $Q$ \emph{avoids} a set $A$ if $A\cap(X_Q\cup Y_Q)=\varnothing$.

We use the semidegree condition through the following simple inequality. Let $F$ be a digraph of order $N$, let $C$ be a set of vertices of $F$, and let $0\le p\le1$ be a function on $V(F)$ which vanishes on $C$. Then, for any two sets $X,Y$ with $X\cap Y\subseteq C$,
\begin{equation}\label{general:capacity}
 |X|+|Y|\le N+|C|-p(V(F))+p(X)+p(Y).
\end{equation}
Indeed, a vertex $v\notin C$ lies in at most one of $X$ and $Y$, so $\ind_X(v)+\ind_Y(v)\le1-p(v)+p(v)(\ind_X(v)+\ind_Y(v))$, while a vertex of $C$ contributes at most two to the left-hand side and one to $N+|C|$.

\begin{lemma}[Deleting vertices of maximum weight]\label{general:max-deletion}
Let $F$ be a digraph of order $N$ with $\delta^0(F)>N/2$. Let $w$ be a deficient fractional $H$-cover of $F$ of minimum total weight, with $0<M=\max w\le1$. Let $C$ be the set of vertices of weight $M$, and $h=|C|$.
Suppose that every tight homomorphism has at most one position mapped to $C$.
Take a fractional homomorphic $H$-tiling of maximum total weight, and delete the position mapped to $C$ from each homomorphism in its support which meets $C$. Let $\rho$ be the resulting measure on partial homomorphisms into $F-C$, and $p$ its multiplicity function.
Then every minimum partial homomorphism $Q$ into $F-C$ satisfies $\Phi_p(Q)>(r-2)h$.
\end{lemma}
\begin{proof}
By Lemma~\ref{general:LP}\ref{LP:slack}, every vertex of $C$ is covered with total multiplicity one, and every homomorphism in the support of the tiling is tight.
Hence $\rho$ has mass $h$, and each partial homomorphism in its support has cost $1-M$.
On the other hand, every partial homomorphism into $F$ has cost at least $1-M$. Indeed, its two completion sets
in $F$ have more than $N/2$ vertices each, so they intersect, and a common vertex, of weight at most $M$, completes it to a homomorphism of $H$, which costs at least one. Therefore, the minimum partial homomorphisms into $F-C$ are exactly the partial homomorphisms into $F-C$ of cost $1-M$.

Suppose, for contradiction, that one of them, $Q$ say, has $\Phi_p(Q)\le(r-2)h$, and let $X$ and $Y$ be its two completion sets in $F$. A vertex of $X\cap Y$ completes $Q$, so it has weight at least $1-(1-M)=M$, and hence
$X\cap Y\subseteq C$.
Extend $p$ by zero on $C$. Since deleting positions only
decreases multiplicities, we have $0\le p\le1$. Moreover,
\[
 p(V(F))=(r-1)h\qquad\text{and}\qquad p(X)+p(Y)\le(r-2)h.
\]
Therefore, \eqref{general:capacity} gives
$2\delta^0(F)\le|X|+|Y|\le N+h-(r-1)h+(r-2)h=N$, a contradiction.
\end{proof}

When a digraph receives no homomorphism from $H$, it contains no large tournament, and we will combine this with the following weighted version of Tur\'an's theorem, which is essentially the Motzkin--Straus theorem.

\begin{lemma}\label{general:weighted-turan}
Let $G$ be a $K_{t+1}$-free graph, and let $p\ge0$ be a function on $V(G)$ with total $P>0$. Then some vertex $v$ with $p(v)>0$ satisfies $p(N_G(v))\le(1-1/t)P$.
\end{lemma}
\begin{proof}
We first show that $\sum_{uv\in E(G)}p(u)p(v)\le(1-1/t)P^2/2$ for every such function $p$.
If $u$ and $v$ are nonadjacent and both have positive value, then, as $p(u)+p(v)$ is kept fixed, this sum is affine in $p(u)$. Thus, moving all the value of one of $u,v$ to the other, chosen appropriately, does not decrease the sum. Repeating this, we reach a function whose support is a clique, which has at most $t$ vertices since $G$ is $K_{t+1}$-free. For such a function, twice the sum is $P^2-\sum_vp(v)^2\le(1-1/t)P^2$, by the Cauchy--Schwarz inequality.
Now, for the original function $p$, the average of $p(N_G(v))$ over the vertices $v$, weighted by $p(v)/P$, is $2\sum_{uv\in E(G)}p(u)p(v)/P\le(1-1/t)P$, which proves the lemma.
\end{proof}

\begin{lemma}
\label{general:tournament-paths}
For each $q\ge1$, every tournament on $q+3$ vertices receives a homomorphism from $D_{q+2,q}$.
\end{lemma}
\begin{proof}
We use induction on $q$. For $q=1$, we need to find $D_{3,1}=Q$ in a tournament on four vertices. Every tournament has a directed Hamilton path, $abcd$ say. If $ad$ is an edge, then this path and the edge $ad$ form a copy of $D_{3,1}$. Otherwise $da$ is an edge, and we look at the edges between $a$ and $c$ and between $b$ and $d$. If $ac$ and $bd$ are edges, then the path $bdac$ and the edge $bc$ form a copy of $D_{3,1}$. Similarly, if $ac$ and $db$ are edges we use $acdb$ and $ab$, if $ca$ and $bd$ are edges we use $cabd$ and $cd$, and if $ca$ and $db$ are edges we use $dbca$ and $da$.

Now let $q\ge2$, let $T$ be a tournament on $q+3$ vertices, and choose a vertex $v$ of $T$, which we take to be a source of $T$ if $T$ has one. By induction, $T-v$ receives a homomorphism from $D_{q+1,q-1}$, and we let $\sigma$ be the image of the common initial vertex of its two paths. Then $\sigma$ has an inneighbour $\sigma'$ in $T$: either $v$ is a source and $\sigma'=v$, or $T$ has no source. Prepending a new initial vertex mapped to $\sigma'$ to both paths gives a homomorphism from $D_{q+2,q}$ to $T$, where the second vertices of the two paths are both mapped to $\sigma$.
\end{proof}

\subsection{Weights in neighbourhoods}\label{sec:prelim-weights}
As every orientation of an even cycle maps onto a $2$-cycle, the neighbourhoods of a vertex of weight zero cannot contain many vertices of small weight, as follows.
For a vertex weighting $w$ of a digraph $F$, we write
\[
 \kappa=\min_{xy\in E(F)}(w(x)+w(y))\qquad\text{and}\qquad
 \chi=\max_{v\in V(F),\,\sigma\in\{+,-\}}\ \min_{u\in N^\sigma(v)}w(u),
\]
so that $\kappa$ is the minimum total weight of an edge, and every vertex has an outneighbour and an inneighbour of weight at most $\chi$.

\begin{lemma}\label{norm391:savings}
Let $H$ be an orientation of an even cycle of order $2S$.
Let $F$ be a digraph of order $N$ with minimum semidegree $d>N/2$,
and let $w$ be a fractional $H$-cover of $F$ with
$\min w=0$ and total weight $W<N/(2S)$. Put $\tau=W/d$, so that $\tau<1/S$.
Then every $2$-cycle $xy$ of $F$ satisfies
$w(x)+w(y)\ge1/S$.
Moreover, let $o$ be a vertex of weight zero, let $\alpha$ and $\alpha'$ be the minimum weights in $N^-(o)$ and $N^+(o)$, and put $m=\min(\alpha,\alpha')$. Then
\[
 \alpha+\alpha'<\tau,\qquad \chi+m<\tau,\qquad m\ge\kappa.
\]
Consequently, $2\kappa<\tau$, $\chi+\kappa<\tau$ and
$S\kappa<1/2$.
\end{lemma}
\begin{proof}
Mapping the two bipartition classes of $H$, each of size $S$, to the two vertices of a $2$-cycle $xy$ gives a homomorphism of cost $S(w(x)+w(y))$, which proves the first assertion.
In particular, every vertex of $N^-(o)\cap N^+(o)$ has weight at least $1/S>\tau$.

Choose sets of $d$ vertices in $N^-(o)$ and in $N^+(o)$, and let $T$ be their intersection, which has size $j>0$ as $2d>N$. Counting the weights in $T$ and in the two other parts of the union of these sets gives
\[
 W\ge j/S+(d-j)(\alpha+\alpha').
\]
If $j=d$, this gives $W\ge d/S>d\tau=W$, a contradiction. Thus $j<d$, and, as $W=d\tau<j/S+(d-j)\tau$, we have $\alpha+\alpha'<\tau$.

Now let $v$ be a vertex, let $\sigma\in\{+,-\}$, and let $\gamma$ be the minimum weight in $N^\sigma(v)$.
As $|N^\sigma(v)|\ge d$, we can choose $d-j$ vertices of $N^\sigma(v)$ outside $T$, and these together with $T$ form a set $A$ of $d$ vertices. The union of the two chosen sets of $d$ vertices around $o$ has $2d-j$ vertices, so at least $d-j$ of them lie outside $A$, and each of these has weight at least $m$. Therefore,
\[
 W\ge j/S+(d-j)(\gamma+m),
\]
and, as before, this gives $\gamma+m<\tau$. Maximising over $v$ and $\sigma$ gives $\chi+m<\tau$.
Finally, $o$ has weight zero, so the edges between $o$ and the vertices attaining $\alpha$ and $\alpha'$ have weights $\alpha$ and $\alpha'$, which gives $\kappa\le m$ and
$2\kappa\le\alpha+\alpha'<\tau$. The remaining assertions follow from these inequalities and $\tau<1/S$.
\end{proof}

\subsection{Homomorphisms into a vertex and an edge}
For the families with span three, we fix an edge $x\to y$, and bound from below the weight of a vertex $z$ which lies in several of the neighbourhoods of $x$ and $y$, by mapping $H$ into $\{x,y,z\}$.

Let $B_*$ be the digraph with vertices $x,y,z$ and edges
$x\to y$, $x\leftrightarrow z$ and $y\leftrightarrow z$.
We say that a random homomorphism from $H$ to $B_*$ is \emph{balanced} if each
vertex of $B_*$ receives $|H|/3$ positions of $H$ in expectation.

\begin{lemma}
\label{compact390:balanced-template}
Let $H$ be an orientation of an even cycle of order $r$ which is not a directed cycle, and whose directed runs all have length at most four. Then there is a balanced random homomorphism from $H$ to $B_*$ which takes three values.
In particular, this applies to $H=+U+V$, where $U$ and $V$ are closed
height walks on $\{0,1,2,3\}$ which start and end at height zero and both visit height three.
\end{lemma}
\begin{proof}
We first construct a homomorphism $f$ from $H$ to $B_*$ whose three fibres have roughly equal size. Let $c=\lfloor r/3\rfloor$. We choose a set $J$ of $c$ positions of $H$, at cyclic distance at least three from each other. We want each path in $H-J$ between two consecutive positions of $J$ to be \emph{alternating}, that is, each of its internal vertices should be a turn of $H$. We then map $J$ to $z$, the sources of the paths of $H-J$ to $x$ and their sinks to $y$. Every edge of $H-J$ goes from a source to a sink, and so maps to $x\to y$, while the other edges map to the $2$-cycles $x\leftrightarrow z$ and $y\leftrightarrow z$. Thus, $f$ is a homomorphism.

To choose $J$, note that consecutive positions of $J$ at distance three leave one edge between them, which is alternating. At distance four they leave a path with two edges, which is alternating if its middle vertex is a turn. At distance five they leave a path with three edges, which is alternating if its two middle vertices are turns.
If $r\equiv0\pmod3$, we take all distances to be three.
If $r\equiv1\pmod3$, we take one distance four, centred at a turn, and all other distances three; a turn exists because $H$ is not a directed cycle.

Suppose that $r\equiv2\pmod3$, so that $r\ge8$. If $H$ has a directed run of length one, its two end vertices are consecutive turns, and we take one distance five, with these turns as the two middle retained positions, and all other distances three.
Otherwise, we find two turns at cyclic distance $L$, where $L\ge4$, $r-L\ge4$ and $L\equiv1\pmod 3$. We then take two distances four, centred at these two turns, and fill the remaining lengths $L-4$ and $r-L-4$, which are nonnegative multiples of three, with distances three. If $H$ has a run of length four, its two end turns are at distance $L=4$, and $r-L\ge4$. Otherwise, all runs have length two or three. Let $n_2$ and $n_3$ be the numbers of runs of these lengths. As $r=2n_2+3n_3$ is even, $n_3$ is even; as the number of runs of a cycle orientation is even, $n_2$ is even too. Since $2n_2+3n_3\equiv2\pmod3$, we have $n_2\equiv1\pmod3$, and so $n_2\ge4$. Take two runs of length two which are consecutive among the runs of length two, with $j$ runs of length three between them. Their outer turns are at distance $L=4+3j$, and the complementary interval contains two further runs of length two, so $r-L\ge4$.

Let $f$ have fibre sizes $(A,A',c)$ at $(x,y,z)$, and put $d=A-A'$. A path of $H-J$ with one or three edges has equally many sources and sinks, while a path with two edges contributes $\pm1$ to $d$. Hence $|d|\le r-3c$.
The two homomorphisms which map the bipartition classes of $H$ to the $2$-cycles $x\leftrightarrow z$ and $y\leftrightarrow z$ have fibre sizes $(r/2,0,r/2)$ and $(0,r/2,r/2)$. With $K=3(r-2c)>0$, choose $f$ with probability $r/K$ and these two homomorphisms with probabilities
\[
 \frac{r-3c-d}K\qquad\text{and}\qquad\frac{r-3c+d}K,
\]
respectively. These probabilities are nonnegative and sum to one, and, since $A=(r-c+d)/2$ and $A'=(r-c-d)/2$, a direct calculation shows that each of $x,y,z$ receives $r/3$ positions in expectation.

Finally, let $H=+U+V$ be as in the statement. A closed walk on $\{0,1,2,3\}$ has directed runs of length at most three, and $U$ and $V$ each start with a positive step and end with a negative step. Thus, the directed run containing a separating positive edge of $+U+V$ consists of this edge and the initial positive run of the following walk, and so it has length at most four. As $H$ is not a directed cycle, the first part of the lemma applies.
\end{proof}

We use Lemma~\ref{compact390:balanced-template} in the following way. Suppose that $\omega\ge0$ is a weighting of a digraph under which every homomorphism of $H$ has cost at least $r/2$, and that $x\to y$ is an edge. If $z$ is a vertex with $x\leftrightarrow z$ and $y\leftrightarrow z$, then $x,y,z$ span a copy of $B_*$, and averaging the cost of the balanced random homomorphism gives $(r/3)\{\omega(x)+\omega(y)+\omega(z)\}\ge r/2$, that is,
\begin{equation}\label{balanced:bound}
 \omega(z)\ge3/2-\omega(x)-\omega(y).
\end{equation}

The vertices which lie in three of the four neighbourhoods of the edge can be bounded similarly, using the directed runs of $H$.
Here and later, for sets $P_1,\ldots,P_\ell$ we write $P_1P_2\cdots P_\ell=P_1\cap\cdots\cap P_\ell$, and, for a set $P$ and a real number $c$, we write $P\ge c$ if $\omega(z)\ge c$ for every $z\in P$, and say that $c$ is a \emph{lower bound} on $P$.

\begin{lemma}
\label{geom368:unrestricted-folds}
Let $H$ be an orientation of an even cycle of order $2S$ which is not a directed cycle,
with $t$ sources and directed runs of lengths $j_1,\ldots,j_{2t}$.
Let $\omega\ge0$ be a weighting of a digraph $D$ under which every homomorphism
from $H$ to $D$ has cost at least $S$. Let $x\to y$ be an edge of $D$, and put
\[
 u=\omega(x),\quad v=\omega(y),\quad k=u+v,
 \quad X=N^-(x),\quad Y=N^+(y),\quad O=N^+(x),\quad I=N^-(y).
\]
Suppose that $A=\sum_h\lfloor j_h/2\rfloor>0$.
Then each of $XYO,XYI,XOI,YOI$ has lower bound
$S(1-k)/A$, and $XYOI$ has lower bound $S(1-k)/A+k/2$.
\end{lemma}
\begin{proof}
Let $B=\sum_h\lfloor(j_h-1)/2\rfloor$ and let $E$ be the number of directed runs of even length. Then $A-B=E$, and, counting the $2t$ turns and the internal positions of the runs, $2t+A+B=2S$. Therefore,
\begin{equation}\label{runs:S}
S=t+B+E/2.
\end{equation}
We write $(xz)^h$ for the sequence $x,z,\ldots,x,z$ with $h$ repetitions of $x,z$.

First let $z\in XOI$, so that $x\to y$, $x\leftrightarrow z$ and $z\to y$ are edges. We send every source of $H$ to $x$ and every sink to $y$, and the positions of a directed run of length $j$, from its source to its sink, to
\[
 (xz)^hxy\ \text{ if }j=2h+1,\qquad (xz)^{h-1}xzy\ \text{ if }j=2h.
\]
Each step of these sequences is one of the available edges, so this is a homomorphism. A run of length $j$ has $\lfloor j/2\rfloor$ internal positions at $z$ and $\lfloor (j-1)/2\rfloor$ at $x$, so the fibre sizes at $(x,y,z)$ are
\[
 (t+B,t,A).
\]

Now let $z\in XYO$, so that $x\to y$, $x\leftrightarrow z$ and $y\to z$ are edges. We \emph{mark} each sink of $H$ which is the end of a directed run of even length, send the marked sinks to $z$, the other sinks to $y$, and all sources to $x$. We send the positions of a run of length $j$ ending at a sink with image $s$ to
\[
 (xz)^hxs\ \text{ if }j=2h+1,\qquad (xz)^{h-1}xyz\ \text{ if }j=2h,
\]
where the second sequence ends at a marked sink. Again every step is an available edge. Each sink is the end of two runs, so the number $m$ of marked sinks satisfies $E/2\le m\le E$, and the fibre sizes are
\[
 (t+B,t+E-m,A-E+m).
\]

In both constructions, the coefficients of $u$ and $v$ are at most $S$, by~\eqref{runs:S} and $m\ge E/2$, while the coefficient $C$ of $\omega(z)$ satisfies $0<C\le A$: for the second construction, $C=B+m$, which is positive since $B=A>0$ if $E=0$, and $m>0$ otherwise.
As every homomorphism has cost at least $S$, we obtain $C\omega(z)\ge S-Sk$, and hence $\omega(z)\ge S(1-k)/A$.
Reversing the sequences above and exchanging the roles of $x$ and $y$ (and marking sources instead of sinks in the second construction) gives homomorphisms which use exactly the edges available for $z\in YOI$ and $z\in XYI$. This proves the other two triple bounds.

Finally, let $z\in XYOI$. Then both the first construction and its reversal are available, with fibre sizes $(t+B,t,A)$ and $(t,t+B,A)$. Averaging their cover inequalities gives
\[
 (t+B/2)k+A\omega(z)\ge S.
\]
Since $t+B/2=S-A/2$ by~\eqref{runs:S}, this gives $\omega(z)\ge S(1-k)/A+k/2$, as required.
\end{proof}
